% Header: General study / work / project template
%%%%%%%%%%%%%%%%%%%%%%%%%%%%%%%%%%%%%%%%%%%%%%%%%%%%%%%%%%%%%%%%%%%

\documentclass[11pt,a4paper]{scrartcl}

% Layout
\usepackage[margin=2.5cm]{geometry}
\usepackage[onehalfspacing]{setspace}

% General packages
\usepackage{graphicx}
\usepackage{enumitem}
\usepackage{tcolorbox}
\tcbuselibrary{breakable}
\usepackage[breaklinks=true,colorlinks=true,linkcolor=blue,urlcolor=blue,citecolor=blue]{hyperref}
\usepackage{amsmath,amsthm,amssymb,mathtools}
\usepackage{braket}
\usepackage{icomma}
\usepackage[T1]{fontenc}
\usepackage[utf8]{inputenc}

% Language: choose one
% \usepackage[ngerman]{babel}
\usepackage[english]{babel}

% Units / tables / floats / references
\usepackage[locale = UK,space-before-unit=true,per-mode = symbol]{siunitx}
\usepackage{booktabs,multirow}
\usepackage{placeins}
\usepackage[natbib,abbreviate=true,doi=false,style=numeric-comp,giveninits=true,sorting=none]{biblatex}
\usepackage{csquotes}

% Optional bibliography
% \addbibresource{references.bib}

% Graphics paths
\graphicspath{{Bilder/} {images/} {figures/} {chapters/}}

% Optional math shortcuts
\newcommand{\R}{\mathbb{R}}
\newcommand{\N}{\mathbb{N}}
\newcommand{\Q}{\mathbb{Q}}
%\newcommand{\set}[1]{\left\{ #1 \right\}}
\newcommand{\abs}[1]{\left| #1 \right|}
\newcommand{\norm}[1]{\left\lVert #1 \right\rVert}
\newcommand{\ol}[1]{\overline{#1}}

% Optional units
\DeclareSIUnit{\dBm}{dBm}

% Box styles
\newtcolorbox{calculationbox}[1][]{
    title=Calculation,
    breakable,
    #1
}

\newtcolorbox{solutionbox}[1][]{
    title=Solution,
    breakable,
    #1
}

\newtcolorbox{answerbox}[1][]{
    title=Answer,
    breakable,
    #1
}

\newtcolorbox{notebox}[1][]{
    title=Notes,
    breakable,
    #1
}

\newtcolorbox{sidecalcbox}[1][]{
    title=Side Calculation,
    breakable,
    #1
}

%%%%%%%%%%%%%%%%%%%%%%%%%%%%%%%%%%%%%%%%%%%%%%%%%%%%%%%%%%%%%%%%%%%
\begin{document}

    \title{Theoretical Physics 2 - Blatt 05}
    \author{Tobias Laser}
    \date{\today}
    \maketitle
    \vfill
    \thispagestyle{empty}

    \pagenumbering{arabic}

    % ================================================================
    \paragraph{13. Bra-Ket Notation} 

    Consider the states $\ket{\psi_n}$ as eigenstates of a Hermitian operator $H$ with 
    \begin{align*}
        H \ket{\psi_n} = \lambda_n \ket{\psi_n}
    \end{align*}

    Assume that the $\ket{\psi_n}$ form a discrete orthonormal basis.
    \begin{enumerate}
        \item The operator $U_{mn}$ is defined as
        \begin{align*}
            U_{mn} = \ket{\psi_m} \bra{\psi_n}.
        \end{align*}

        Compute the adjoint operator $U^\dagger_{mn}$, the commutator $[U_{mn}, H]$, and show that 
        \begin{align*}
            U_{mn} U_{pq} = \delta U_{mq}
        \end{align*}

        (kronecker delta).

        \begin{calculationbox}
            \begin{calculationbox}
                $U^\dagger_{mn} = (\ket{\psi_m} \bra{\psi_n})^\dagger = \bra{\psi_n}^\dagger \ket{\psi_m}^\dagger = \ket{\psi_n} \bra{\psi_m} = U_{nm}$
            \end{calculationbox}

            \begin{calculationbox}
                $U_{mn} U_{pq} = \ket{\psi_m} \braket{\psi_n | \psi_p} \bra{\psi_q} = \ket{\psi_m} \delta_{np} \bra{\psi_q} = \delta_{np} \ket{\psi_m} \bra{\psi_q} = \delta_{np} U_{mq}$
            \end{calculationbox}
            
        \end{calculationbox}

        \item Show that for an arbitrary operator A and 
        \begin{align*}
            a_{mn} = \braket{\psi_m | A | \psi_n},
        \end{align*}

        the following holds:
        \begin{align*}
            A = \sum\limits_{m,n} a_{mn} U_{mn}
        \end{align*}

        \begin{calculationbox}
            Lets define $B := \sum\limits_{m,n} a_{mn} U_{mn}$ and let $B$ operate on an arbitrary basis vector $\ket{\psi_k}$:
            
            $B \ket{\psi_k} 
            = \sum\limits_{m,n} a_{mn} \ket{\psi_m} \braket{\psi_n | \psi_k} 
            = \sum\limits_{m,n} a_{mn} \ket{\psi_m} \delta_{nk} 
            = \sum\limits_m a_{mk} \ket{\psi_m} 
            = \sum\limits_n \braket{\psi_m | A | \psi_k} \ket{\psi_m}
            = \sum\limits_n \ket{\psi_m}\braket{\psi_m | A | \psi_k}
            = A \ket{\psi_k}$
            \begin{calculationbox}
                $\ket{\phi} = \sum\limits_m \ket{\psi_m} \braket{\psi_m |\phi}$
            \end{calculationbox}

            Therefore $a = B = \sum\limits_{m,n} a_{mn} U_{mn}$
        \end{calculationbox}

        or other way of calculation

        \begin{calculationbox}
            Completeness relation $I = \sum\limits_m \ket{\psi_m} \bra{\psi_m}$

            $A 
            = I A I 
            = (\sum\limits_m \ket{\psi_m} \bra{\psi_m})A(\sum\limits_n \ket{\psi_n} \bra{\psi_n}) 
            = \sum\limits_{m,n} \ket{\psi_m} \braket{\psi_m |A | \psi_n} \bra{\psi_n} 
            = \sum\limits_{m,n} \ket{\psi_m} a_{mn} \bra{\psi_n} 
            = \sum\limits_{m,n} a_{mn} \ket{\psi_m | \psi_n} 
            = \sum\limits_{mn} a_{mn} U_{mn}$
        \end{calculationbox}

        \item Consider the following states:
        \begin{align*}
            \ket{\phi} &= \frac{1}{\sqrt{2}} \ket{\psi_1} + \frac{i}{2} \ket{\psi_2} + \frac{1}{2} \ket{\psi_3},\\
            \ket{X} &= \frac{1}{4} \ket{\psi_1} + \frac{1}{3} \ket{\psi_3}. 
        \end{align*}

        Are these states normalized? 
        \begin{calculationbox}
            $\braket{\phi | \phi} \stackrel{!}{=} 1$
            \begin{calculationbox}
                $\braket{\phi | \phi}
                = \abs{\frac{1}{\sqrt{2}}}^2 + \abs{\frac{i}{2}}^2 + \abs{\frac{1}{2}}^2
                = \frac{1}{2} + \frac{1}{4} + \frac{1}{4}
                = 1
                $
            \end{calculationbox}

            $\braket{X | X} \stackrel{!}{=} 1$
            \begin{calculationbox}
                $\braket{X | X}
                = \abs{\frac{1}{4}}^2 + \abs{\frac{1}{3}}^2
                = \frac{1}{16} + \frac{1}{9} \neq 1
                $
            \end{calculationbox}

        \end{calculationbox}

        Are they orthogonal?
        \begin{calculationbox}
            $\braket{\phi | X} \stackrel{!}{=} 0$
            \begin{calculationbox}
                $\braket{\phi | X}
                = \ket{\phi}^\dagger \ket{X}
                = (\frac{1}{\sqrt{2}} \ket{\psi_1} + \frac{i}{2} \ket{\psi_2} + \frac{1}{2} \ket{\psi_3})^\dagger (\frac{1}{4} \ket{\psi_1} + \frac{1}{3} \ket{\psi_3})$

                $= (\frac{1}{\sqrt{2}} \bra{\psi_1} - \frac{i}{2} \bra{\psi_2} + \frac{1}{2} \bra{\psi_3}) (\frac{1}{4} \ket{\psi_1} + \frac{1}{3} \ket{\psi_3})
                =\frac{1}{\sqrt{2}} \frac{1}{4} + \frac{1}{2}\frac{1}{3} > 0 
                $
            \end{calculationbox}

        \end{calculationbox}

        Compute the projectors onto $\ket{\phi}$ and $\ket{X}$, and show that the projectors are self-adjoint operators.
        \begin{calculationbox}
            Let $\ket{\psi_1} := \begin{pmatrix}
                1\\0\\0
            \end{pmatrix}$, 
            $\ket{\psi_2} := \begin{pmatrix}
                0\\1\\0
            \end{pmatrix}$,
            $\ket{\psi_3} = \begin{pmatrix}
                0\\0\\1
            \end{pmatrix}$

            \begin{calculationbox}
                $\ket{\phi} = \begin{pmatrix}
                    \frac{1}{\sqrt{2}} \\ \frac{i}{2}\\ \frac{1}{2}
                \end{pmatrix}$, 
                $\bra{\phi} = \begin{pmatrix}
                    \frac{1}{\sqrt{2}} & -\frac{i}{2} & \frac{1}{2}
                \end{pmatrix}$

                $P_\phi 
                = \ket{\phi} \bra{\phi}
                = \begin{pmatrix}
                    \frac{1}{2} & -\frac{i}{2\sqrt{2}} & \frac{1}{2\sqrt{2}}\\
                    \frac{i}{2\sqrt{2}} & \frac{1}{4} & \frac{i}{4}\\
                    \frac{1}{2\sqrt{2}} & -\frac{i}{4} & \frac{1}{4}
                \end{pmatrix}
                $
            \end{calculationbox}

            \begin{calculationbox}
                $\ket{\phi} = \begin{pmatrix}
                    \frac{1}{4} \\ 0 \\ \frac{1}{3}
                \end{pmatrix}$, 
                $\bra{\phi} = \begin{pmatrix}
                    \frac{1}{4} & 0 & \frac{1}{3}
                \end{pmatrix}$

                $P_\phi 
                = \ket{\phi} \bra{\phi}
                = \begin{pmatrix}
                    \frac{1}{16} & 0 & \frac{1}{12}\\
                    0 & 0 & 0 \\
                    \frac{1}{12} & 0 & \frac{1}{9}
                \end{pmatrix}
                $
                
            \end{calculationbox}
        \end{calculationbox}

        \item Let $\ket{\phi}$ and $\ket{X}$ be two arbitrary vectors in the state space. Under what condition is 
        \begin{align*}
            K = \ket{\phi} \bra{X}
        \end{align*}
        Hermitian?

        \begin{calculationbox}
            Condition $K^\dagger \stackrel{!}{=} K$
            \begin{calculationbox}
                $K^\dagger = \ket{X} \bra{\phi}$
            \end{calculationbox}
            That does not hold for all $\ket{\phi}, \ket{X}$. That can only be true if they are linear dependent $\ket{X} = \lambda \ket{\phi}$

            \begin{calculationbox}
                $K = \ket{\phi} \bra{X} = \ket{\phi} (\lambda \ket{\phi})^\dagger = \ket{\phi} \lambda^\star \bar{\phi} = \lambda^\star \ket{\phi} \bar{\phi}$
            \end{calculationbox}
            $\implies \lambda = \lambda^\star$
        \end{calculationbox}

        Under what condition is $K$ a projector? 
        
        \begin{calculationbox}
            $K^2 \stackrel{!}{=}$
            \begin{calculationbox}
                $K^2
                = (\ket{\phi} \bra{X}) (\ket{\phi} \bra{X})
                = \ket{\phi} \braket{X | \phi} \bra{X}
                = \braket{ X | \phi} \ket{\phi} \bra{X}
                = \braket{X | \phi} K
                $
            \end{calculationbox}
            $\implies$
            \begin{calculationbox}
                Trivial for $K = 0 $

                $\braket{X | \phi} K = K \implies \braket{X | \phi} = 1$
            \end{calculationbox}
        \end{calculationbox}

        Show that $K$ can always be written as 
        \begin{align*}
            K = \alpha P_1 P_2
        \end{align*}
        where $P_i$ are projectors and $\alpha \in \mathbb{C}$.

        \begin{calculationbox}
            Start with projections in the direction of $\ket{\phi}$ and $\ket{X}$
            
            \begin{calculationbox}
                For $ket{\phi} \neq 0$ and $\ket{X} \neq 0$

                we get the normalized

                $\ket{\hat{\phi}} = \frac{\ket{\phi}}{\norm{\phi}}$ and 
                $\ket{\hat{X}} = \frac{\ket{X}}{\norm{X}}$
            \end{calculationbox}

            Let $P_1 = \ket{\hat{\phi}} \bra{\hat{\phi}}$,
            $P_2 = \ket{\hat{X}} \bra{\hat{X}}$
            and calculate:
            \begin{calculationbox}

                $P_1 P_2
                = \ket{\hat{\phi}} \braket{\hat{\phi} | \hat{X}} \bra{\hat{X}}
                = \braket{\hat{\psi} | \hat{X}} \ket{\hat{\phi}} \bra{\hat{X}}.
                $

                $\implies$

                $\ket{\hat{\phi}} \bra{\hat{X}}
                = \frac{1}{\norm{\phi}\norm{X}} \ket{\phi} \bra{X}
                $

                $P_1 P_2 = \frac{\braket{\hat{\phi} | \hat{X}}}{\norm{\phi} \norm{X}} \ket{\phi} \bra{X}$
            \end{calculationbox}
            With this $K = \ket{\phi} \bra{X} 
            = \frac{\norm{\phi} \norm{X}}{\braket{\hat{\phi} | \hat{X}}} P_1 P_2
            = \frac{\norm{\phi}^2 \norm{X}^2}{\braket{\phi | X}} P_1 P_2
            = \alpha P_1 P_2$
        \end{calculationbox}

        \item Let $H \ket{\phi_n} = \lambda_n \ket{\phi_n}$ and assume that the eigenvalues $\lambda$  are non-degenerate. Show that the so-calles \textbf{null operator}: 
        \begin{align*}
            \Pi_0 = \prod\limits_n (H - \lambda_n)
        \end{align*}
        satisfies
        \begin{align*}
            \Pi_0 \ket{\phi} = 0
        \end{align*}
        for all vectors $\ket{\phi}$.

        \begin{calculationbox}
            Want to show, that when $\ket{\phi_n}$ is as basis, that $\Pi_0 = 0$ on any basis vectors. 
            \begin{calculationbox}
                Take an arbitrary eigenvector $\ket{\phi_m}$ and lets apply $H-\lambda_n$ on it:

                $(H - \lambda_n) \ket{\phi_m} = (\lambda_m - \lambda_n)\ket{\phi_m}$,
            \end{calculationbox}
            
            Lets look at the null operator
            \begin{calculationbox}
                $\Pi_0\ket{\phi_m} 
                = \prod\limits_n (H - \lambda_n) \ket{\phi_m}
                = \prod\limits_n (\lambda_m - \lambda_n) \ket{\phi_m}
                = ... \cdot (\lambda_m - lambda_m) \cdot ...
                = 0$
            \end{calculationbox}
            and applied to any vector $\ket{\phi} = \sum\limits_m c_m \ket{\phi_m}$
            \begin{calculationbox}
                $\Pi_0 \ket{\phi}
                = \Pi_0 \sum\limits_m c_m \ket{\phi_m}
                = \sum\limits_m c_m \Pi_0 \ket{text}
                =\sum\limits_m c_m \cdot 0
                = 0
                $
            \end{calculationbox}

        \end{calculationbox}

        How does the operator
        \begin{align*}
            \Pi_{n \neq m} \left( \frac{H - \lambda_n}{\lambda_m - \lambda_n} \right)
        \end{align*}
        act?

        \begin{calculationbox}
            Let $P_m := \prod\limits_{n \neq m} \left( \frac{H - \lambda_n}{\lambda_m - \lambda_n} \right)$

            Lets apply $P_m$ an an arbitrary eigenvector $\ket{\phi_k}$
            \begin{calculationbox}
                $P_m \ket{\phi}$

                Case 1: $k = m$
                \begin{calculationbox}
                    $P_m \ket{\phi_m} = \prod\limits_{n \neq m} \left( \frac{\lambda_m - \lambda_n}{\lambda_m - \lambda_n} \ket{\phi_m} \right)$

                    $\implies$

                    $\P_m \ket{\phi_m} = 1 \ket{\phi_m}$
                \end{calculationbox}

                Case 2: $k \neq m$
                \begin{calculationbox}
                    $\lambda_k - \lambda_k = 0$
                    Therefore $P_m \ket{\phi_k} = 0$
                \end{calculationbox}
            \end{calculationbox}
            
            For any arbitrary vector $\ket{\phi} = \sum\limits_k c_k \ket{\phi_k}$ the operator takes only the m-th component of the vector.
            
            $P_m \ket{\phi}
            =\sum\limits_k c_k P_m \ket{\phi_k} \sum_k \delta_{mk} \ket{\phi_k}
            = c_m \ket{\phi_m}
            $
        \end{calculationbox}

        Can this operator be expressed in another way?
        \begin{calculationbox}
            $P_m = \ket{\phi_m} \bra{\phi_m}$

            $\ket{\phi_m} \bra{\phi_m} \ket{\phi_k} = \delta_{mk} \ket{\phi_m}$

            Same effect, because it takes the k-th component of the vector $\ket{\phi}$

        \end{calculationbox}

        \item  Compute the matrix elements of the momentum operator $\hat{p}$ in the position representation,
        \begin{align*}
            \braket{x' | \hat{p} | x},
        \end{align*}
        using the momentum representation
        \begin{align}
            \braket{p'|p} = \delta(p' - p),
        \end{align}
        and the known wavefunction
        \begin{align*}
            \braket{x | p} = \frac{\exp{\frac{ipx}{\hbar}}}{\sqrt{2\pi \hbar}}
        \end{align*}

        \begin{calculationbox}
            $1 = \int dp \ket{p}\bra{p}$

            $\braket{x'|1 \cdot \hat{p} \dot 1|x} = \int dp' \int dp \braket{x'|p'} \braket{p'|\hat{p} | p} \braket{p|x}$

            \begin{calculationbox}
                Eigenstade/ eigenvector $\ket{p}$

                $\hat{p}\ket{p} = p \ket{p} \implies \braket{p'| \hat{p} | p} = p\braket{p'| p} = p \delta(p' - p)$
            \end{calculationbox}

            Plug in
            $\braket{x' | \hat{p} | x} = \int dp' \int dp \braket{x'|p'} p \delta(p'-p) \braket{p|x} = \int dp \braket{x' | p} p \braket{p | x}$

            plug in wavefunction and adjungated wavefunction ($\braket{p | x} = \frac{\exp{\frac{-ipx}{\hbar}}}{\sqrt{2\pi \hbar}}$)

            $\braket{x' | \hat{p} | x} 
            = \int dp \frac{\exp{\frac{ipx'}{\hbar}}}{\sqrt{2\pi \hbar}} p \frac{\exp{\frac{-ipx}{\hbar}}}{\sqrt{2\pi \hbar}}
            = \frac{1}{2 \pi \hbar} \int dp p \exp{\frac{ip (x'-x)}{\hbar}}
            $

            Using Identity: $\delta(x' - x) = \frac{1}{2\pi\hbar} \int dp \exp{\frac{i p (x' - x)}{\hbar}}$

            the derivation of the Identity $\frac{\partial}{\partial x'} \delta(x' - x) = \frac{1}{2\pi\hbar} \int dp \frac{ip}{\hbar} \exp{\frac{i p (x' - x)}{\hbar}}$

            And therefore $\frac{1}{2\pi \hbar} \in dp p \exp{\frac{i p (x' - x)}{\hbar}} = - i \hbar \frac{partial}{\partial x'} \delta(x' - x)$

            Result $\braket{x' | \hat{p} | x} = - i \hbar \frac{\partial}{\partial x'} \delta(x' - x)$

        \end{calculationbox}

    \end{enumerate}

\end{document}